\section{From Marital Assortment to Historical Explanations}
\label{sec:historical-application}
Chapter~18 moves from family resemblance to a claim about how marital
assortment changes the supply of exceptional abilities. This application
requires a separate chain of evidence: the genetic nature of the trait,
the degree and history of assortment, the change in its distribution,
and the connection from that distribution to innovation. A successful
fit in Part~I would not establish all these links.

\subsection{A Classical Variance Result and a Numerical Correction}
An increase in equilibrium genetic variance under positive assortment is
a classical result of quantitative genetics, originating in Fisher's
framework and developed in subsequent work \citep{Fisher1918,Nagylaki1982}.
Clark's contribution here is the proposed historical application, whose
premises require separate evaluation. In the chapter's simplified model,
with fixed segregation variance $S$,
\[
 V_{t+1}=\frac{1+m}{2}V_t+S,
 \qquad V(m)=\frac{2S}{1-m}.
\]
Thus $V(.6)/V(0)=2.5$, not the fivefold increase stated on p.~274
\citep{Clark2026}. The fivefold result is correct for $m=.8$, the value
used in \citet[pp.~3--4]{ClarkCummins2022}. The arithmetic correction is
separate from whether fixed segregation variance is an appropriate
approximation for the historical comparison.

\subsection{Marriage Practices Do Not Directly Measure Genetic Assortment}
Observed spouse resemblance, a latent genetic correlation, and genetic
relatedness through shared ancestors are different quantities. Comparing
them across societies requires a common trait definition and an account
of measurement, partner selection, and population history. Section~\ref{sec:genomic-evidence}
explains how family-based PGIs could supply evidence about one of these
links; an inference from marriage institutions alone does not.

The cousin-marriage calculation makes the difficulty concrete. Clark
assumes random choice within the cousin pool and substitutes the cousin
correlation into a fixed-point equation for spouse genetic resemblance
\citep[pp.~278--280]{Clark2026}. Random choice with respect to the relevant
trait is an assumption, not a consequence of marrying a cousin. Moreover,
the cousin formula derived for pedigrees with independent outside mates
cannot simply be carried over to recurrent consanguineous mating: repeated
shared ancestors change the paths of relationship. This is a requirement
for a new derivation, not a claim that a corrected equilibrium has already
been established here. The calculation does not by itself demonstrate
weaker genetic assortment outside England.

\subsection{Variance, Exceptional Outcomes, and Historical Causation}
Even a justified increase in variance would not determine the number of
innovators without a model relating measured ability to exceptional
achievement. Tail probabilities depend on distributional assumptions,
and achievement also depends on opportunities to develop and apply a
trait. A comparison of equilibrium distributions is not evidence that
historical societies occupied those equilibria or that their differences
caused industrialization. These are additional empirical claims, which
should be evaluated separately from the classical variance result.
